\section{Prior Literature and the Algebraic Target}\label{sec:mergeable-sketches}

The framework is easiest to read as a bridge from five older ideas to one
learned object. Database systems already distinguish aggregates whose
fixed-size intermediate state can be combined before finalization
\citep{GrayEtAl1997DataCube}. Streaming theory formalizes this as
mergeable summaries \citep{Agarwal2013MergeableTODS}; distributed
streaming work studies the symmetric unordered case
\citep{FeldmanEtAl2008MUD}; and functional programming gives the ordered
analogue through list, or free-monoid, homomorphisms
\citep{Gibbons1996ThirdHomomorphism}. Finally, the sufficient-statistic
tradition explains why a compressed state is useful only relative to a
downstream target \citep{Fisher1922,Blackwell1953}. C-TreePO takes this
lineage as its algebraic target: learn a bounded state, merge it locally,
and audit whether the task readout is preserved.

\subsection{Mergeable Summaries in Data Systems}

A \emph{mergeable summary} is a compact data structure equipped with a
binary merge operation that behaves well under hierarchical aggregation.
\citet{Agarwal2013MergeableTODS} give the definition that modern streaming
literature has inherited: a summary scheme of size $k$ for a query $Q$
(cardinality, frequency, quantile, membership, \dots) is \emph{mergeable}
if for any two streams $u, v$ there is a merge operator $\oplus$ such that
$\mathrm{sketch}(u \concat v)$ and
$\mathrm{sketch}(u) \oplus \mathrm{sketch}(v)$ both answer $Q(u \concat v)$
equally well, with sketch size bounded by $k$ after merging. The appeal is
operational: a MapReduce, Dataflow, or Spark job can compute a sketch per
shard, merge shards pairwise up a tree, and read the query off the root,
without ever materializing the full stream. The same state-composition
pattern appears in massive unordered distributed aggregation and in database
``algebraic'' aggregates, where fixed-size intermediate state is combined
before finalization \citep{FeldmanEtAl2008MUD,GrayEtAl1997DataCube}.

The mergeable-summary literature is deep. HyperLogLog
\citep{FlajoletEtAl2007,HeuleEtAl2013} estimates distinct counts in a few
kilobytes regardless of stream length. Count-Min
\citep{CormodeMuthukrishnan2005} answers frequency queries under a space
budget. Bloom filters, Theta/KMV sketches, t-digest, KLL, and GK quantile
sketches are part of the same systems lane, with different exact or
approximate merge guarantees. The common thread is that the state,
merge, and readout are engineered before training, so correctness is
usually proved from sketch-specific algebra rather than inferred from
runtime audits. A machine-checked layer of this literature (which
sketches carry Lean-backed merge laws, and which are wrapped only as
empirical references) is inventoried in
Appendix~\ref{app:theory-details}. HLL is
the empirical anchor below, not the definition of the classical class.

\subsection{The Algebraic View}

Abstractly, a mergeable sketch for a query family indexed by $q \in \mathcal{Q}$
is a state-level triple $(\mathrm{encode}, \oplus, \mathrm{query})$ with
\[
\mathrm{encode} : \Strings \to \mathcal{S}, \quad
\oplus : \mathcal{S} \times \mathcal{S} \to \mathcal{S}, \quad
\mathrm{query} : \mathcal{S} \times \mathcal{Q} \to \mathcal{R},
\]
and correctness
\[
\mathrm{query}\big(\mathrm{encode}(u \concat v),\, q\big)
\;\approx\;
\mathrm{query}\big(\mathrm{encode}(u) \oplus \mathrm{encode}(v),\, q\big)
\]
for every pair of shards $u, v$ and every query $q$. The merge composes
\emph{states}, not scalar query answers. Distinct count is the canonical
warning: $|A \cup B|$ is not determined by $|A|$ and $|B|$ alone, while an
HLL register vector can still be merged because it carries overlap-relevant
state. Associativity of $\oplus$ is the schedule-invariance condition: any
binary aggregation tree reaches the same state, or at least the same readout,
so the merge schedule can be chosen for execution convenience without
changing the answer. Commutativity is specific to unordered or symmetric
inputs; it lets shards be merged without tracking their original order.
Idempotence is specific to set- or semilattice-like unions such as HLL or
Bloom-filter-style membership sketches, where it makes re-merging an
already-absorbed shard harmless; frequency and quantile summaries lack it.

For ordered text, the corresponding algebraic ideal is a list, or free-monoid,
homomorphism: $h(u \concat v) = h(u) \odot h(v)$ for an associative operator
$\odot$ on summary states \citep{Gibbons1996ThirdHomomorphism}. The operator
$\odot$ need not be commutative, because word and paragraph order can matter.
C-TreePO is an audited learned approximation to this state-level discipline:
the summaries are learned states, and the local laws test whether those states
still support the task readout.

HyperLogLog makes the story concrete: a sketch of precision $p$ keeps
$m = 2^p$ registers, encodes each item by a hash-indexed
leading-zero count, merges by pointwise max (associative, commutative,
and idempotent, so every tree-shaped reduction yields the same register
vector), and reads cardinality off a harmonic-mean formula with
relative standard error $1.04/\sqrt{m}$ regardless of stream content
\citep{FlajoletEtAl2007}. This is the pattern C-TreePO generalizes: a
bounded-size state, an associative merge, and a query that reads out a
task value. Because encode, merge, and query are deterministic and
closed-form, a single algebraic proof covers every realized shard and
no runtime audit is needed. Appendix~\ref{app:classical-parity} gives
the full mechanics and the parity experiments.

\subsection{The C-TreePO Correspondence}

C-TreePO lifts this recipe to settings where no hand-designed sketch is
available. The leaf state is a learned summary $g(b)$ of a raw text block;
the merge operator is the same learned summarizer applied to a
concatenation $g(s_{u_L} \concat s_{u_R})$; the query is an oracle value
$\fstar(x)$ that a domain expert would assign to the raw document.
Table~\ref{tab:sketch-mapping} makes the mapping explicit.

\begin{table}[htbp]
    \centering
    \small
    \begin{tabular}{p{0.46\linewidth}p{0.46\linewidth}}
        \toprule
        Classical mergeable sketch & C-TreePO semantic compression \\
        \midrule
        Data shard $d_i$ & Leaf text block $b_i$ \\
        Bounded sketch state $S_B(d_i)$ & Leaf summary $s_i = g(b_i)$ \\
        Merge operator $\oplus_B$ & Learned merge $g(s_{u_L}\concat s_{u_R})$ \\
        Query (e.g., distinct count, frequency) & Oracle query $\fstar(x)$ (task value) \\
        Memory budget $B$ controls what survives merges & Summary budget controls interaction-bearing content retention \\
        Correctness by algebraic proof (no runtime check) & Correctness by probabilistic audit (Section~\ref{sec:estimation}) \\
        Error = sketch estimation variance & Error = $L \cdot \Delta_R$ + estimation envelopes (Theorem~\ref{thm:e2e}) \\
        \bottomrule
    \end{tabular}
    \caption{Mapping from a classical mergeable sketch to the C-TreePO setting. In classical sketches, the merge operator is hand-designed and algebraically exact; in C-TreePO, it is learned and certified by audit.}
    \label{tab:sketch-mapping}
\end{table}

Two differences matter. First, the merge $g$ is \emph{learned} rather than
hand-designed: a neural network operating on tokens. Its correctness has no
sketch-specific algebra to lean on, so Section~\ref{sec:estimation} replaces the
algebraic proof by a probabilistic audit that samples realized merges and
estimates how often they preserve the oracle value. In the optimization
language of Section~\ref{sec:theorems}, the local-law weights are the
projection pressure that moves this learned operator toward the same
theorem-eligible subspace that a classical sketch occupies by construction.
Second, the readout is a \emph{task value}, not necessarily a discrete
aggregate. For HyperLogLog the readout is a distinct-count estimate; for a
political-science application it might be a manifesto's RILE score
\citep{BenoitEtAl2025,ManifestoProject2025a}; for the synthetic task of
Appendix~\ref{sec:markov-interlude} it is a count of color-regime transitions.
The common discipline is state composition plus readout preservation. Whether
a compact state exists is task-specific and is exactly what the audit probes
when no hand-designed sketch is available.

Classical full mergeable summaries nest in C-TreePO through relational
state-level C-Trees. Randomized mergeable summaries nest in the same
formulation in probability: root readout is correct on the event that the
merged root state is valid. This statement preserves the union
formalization: shards combine as concatenated datasets, the approximation
parameter $\varepsilon$ stays fixed in the size/error profile, and
sample-size or training-data requirements are separate concentration
side conditions rather than scalar child-query merge laws.

\subsection{The Generalization Claim}

There are two useful limiting cases. The first is a strict oracle
homomorphism, where the task value itself has a merge operation. This is
mathematically clean but too strong for many classical sketches: scalar
distinct count, for example, cannot merge $|A|$ and $|B|$ into $|A\cup B|$
without overlap information. The second is the actual mergeable-sketch limit,
where the \emph{state} carries enough information for a bounded readout. The
following proposition keeps these two cases separate.

\begin{proposition}[Strict and State-Level Mergeable Limits]\label{prop:mergeable-reduction}
Under Assumption~\ref{ass:context} (context-compatible oracle), suppose $g$ is deterministic and satisfies global versions of the preservation laws:
\[
\text{(A1)}\ \forall z,\ g(z) \sim z,\qquad
\text{(A2)}\ \forall u,v,\ u \concat v \sim g(g(u)\concat g(v)),
\]
Define the induced summary merge $\oplus_g(s,t):=g(s\concat t)$. Then:
\begin{enumerate}[nosep]
\item \emph{Strict oracle homomorphism.} If there exists
$m_{\Y}:\Y\times\Y\to\Y$ such that
$m_{\Y}(\fstar(u),\fstar(v))=\fstar(u\concat v)$ for all $u,v$, then
$\fstar$ is a homomorphic readout of the induced state: merging summaries and
then reading out gives the same oracle value as reading the raw concatenation,
and $\oplus_g$ is associative modulo $\sim$.
\item \emph{Classical state-level reduction.} More generally, let
$(\mathrm{encode},\oplus_B,\mathrm{query})$ be a bounded mergeable-summary
scheme for a query family $Q_q$ in the sense of
\citet{Agarwal2013MergeableTODS}. If $g$ serializes the sketch state
($g(x)=\mathrm{encode}(x)$ on raw shards and
$g(s\concat t)=s\oplus_B t$ on serialized states) and the C-TreePO readout is
$f(s,q)=\mathrm{query}(s,q)$, then C-TreePO is exactly that classical
mergeable summary under the induced merge $\oplus_g$.
\end{enumerate}
\end{proposition}

\noindent\textit{Proof intuition.} The strict case chains A1 and A2 through
context compatibility to recover homomorphic readout and associativity modulo
the oracle. The classical case is state-level: bounded sketch states compose
by $\oplus_B$, and the readout is applied only after the state has been
merged. In the randomized Agarwal setting, the same conclusion holds with the
same success-probability reading: conditional on the root state being valid for the
merged union stream, the root readout equals the target value. Full proof in
Appendix~\ref{app:proof-mergeable}.

Appendix~\ref{sec:hll-parity} exhibits this collapse numerically on one
representative mergeable sketch, HyperLogLog. Running a classical HLL
(native plus the Apache DataSketches reference) through the unified
\texttt{fit()} framework with $n_{\mathrm{leaves}} \in \{1, 2, 4, 8,
16\}$ yields byte-identical register states between TreePO-reduced and
flat pipelines for the register-based implementation, and
estimate-equivalence within the HLL relative standard error for the
DataSketches reference. The same pipeline can swap its oracle $\fstar$
from analytic cardinality to the
HLL's own scoring head, and under that oracle a fully end-to-end learned
$g$ and $f$ track the classical sketch within a small multiple of the RSE
floor. In the projection language, the
classical rows start inside the exact local-law subspace, while the
learned rows test different representable subspaces and readouts for how
closely training can reach that target. The formal claim extends beyond
HLL: any sketch that supplies the
same \texttt{SketchLeafPreserving}, \texttt{SketchMergeCompatible}, and
\texttt{SketchSummaryCompatible} witnesses enters the exact local-law
subspace through
\texttt{mergeableSketchSummaryClass\_subset\_exactLocalLawSubspace}.
A companion benchmark suite, \texttt{treepo-bench classical-sketches},
repeats the tree-reduction comparison for the broader official
DataSketches families through the same \texttt{fit()} interface
(Appendix~\ref{app:treepo-software}).

Most interesting oracles lie beyond register-like queries. In the Markov example of Appendix~\ref{sec:markov-interlude}, the oracle counts boundary transitions between adjacent colored regions---a cross-span interaction that no commutative register-wise merge can recover from leaves alone. The merge operator has to carry boundary information forward, and that is what motivates the Fourier neural operator in Appendix~\ref{sec:fno-primer}. Section~\ref{sec:framework} develops the local laws, Section~\ref{sec:estimation} the audit, and Section~\ref{sec:theorems} the preference-learning consequences that make the learned version as trustworthy as the classical one.
